% =====================================================================
% CV - Jorge Pizarro Callejas (English version)
% ATS-optimized LaTeX template
% Compile with: pdflatex CV_Jorge_Pizarro_Callejas_EN.tex
%
% ATS-oriented design decisions:
%   - Single column, no tables, no floating boxes, no icons
%   - Standard font (Times via mathptmx) with T1 encoding
%   - Glyph-to-Unicode mapping (glyphtounicode) so accented characters
%     extract correctly from the PDF
%   - Conventional section headings
%   - Simple bullets (itemize), left-aligned text
%   - Links shown as plain text (no colored boxes)
% Always verify with:  pdftotext CV_Jorge_Pizarro_Callejas_EN.pdf -
% =====================================================================
\documentclass[10pt,letterpaper]{article}

% --- Encoding and fonts ----------------------------------------------
\usepackage[utf8]{inputenc}
\usepackage[T1]{fontenc}
\usepackage{mathptmx}
\usepackage{microtype}

% --- Reliable text extraction ----------------------------------------
\input{glyphtounicode}
\pdfgentounicode=1

% --- Layout ----------------------------------------------------------
\usepackage[top=1.5cm,bottom=1.5cm,left=1.8cm,right=1.8cm]{geometry}
\usepackage{ragged2e}
\usepackage{enumitem}
\usepackage{titlesec}
\usepackage{xcolor}
\usepackage[hidelinks,pdftitle={CV - Jorge Pizarro Callejas},pdfauthor={Jorge Pizarro Callejas}]{hyperref}

\RaggedRight
\setlength{\parindent}{0pt}
\setlength{\parskip}{0pt}
\pagestyle{empty}
\setlength{\emergencystretch}{3em}

% --- Section style ---------------------------------------------------
\titleformat{\section}{\large\bfseries}{}{0pt}{\MakeUppercase}[\vspace{-0.5em}\rule{\linewidth}{0.4pt}]
\titlespacing*{\section}{0pt}{10pt}{5pt}

% --- Lists -----------------------------------------------------------
\setlist[itemize]{leftmargin=1.2em,itemsep=2pt,topsep=2pt,parsep=0pt,partopsep=0pt}

% --- Reusable macros -------------------------------------------------
% \cvjob{Title}{Company, city}{Dates}
\newcommand{\cvjob}[3]{%
  \vspace{5pt}%
  \textbf{#1} | #2 \hfill #3\par\vspace{1pt}}

% \cvskill{Category}{List of skills}
\newcommand{\cvskill}[2]{\item \textbf{#1:} #2}

% Skill separator (plain text, extractable)
\newcommand{\sep}{\ \textbullet\ }

\begin{document}

% =====================================================================
% HEADER
% =====================================================================
{\LARGE\bfseries Jorge Pizarro Callejas}\par\vspace{3pt}
{\large\bfseries Software Engineer | DevOps \& Cloud (AWS) | Linux}\par\vspace{3pt}
Quilpué, Chile | (+56) 976247629 | \href{mailto:jpizarrocallejas@gmail.com}{jpizarrocallejas@gmail.com} | \href{https://linkedin.com/in/jorgepizarrocallejas}{linkedin.com/in/jorgepizarrocallejas} | \href{https://github.com/jorgicio}{github.com/jorgicio}\par

% =====================================================================
\section{Professional Summary}
% =====================================================================
% TODO: specify which package repository you maintain (the DOCX had
% "[confirm which one]"). Example: "I maintain a package repository on AUR".
Software Engineer (B.Eng. in Computer Science, UTFSM) with a web and software development career since 2011 and 3+ years building and operating services on AWS (Lambda, CloudFormation) with CI/CD on Jenkins and GitLab. At Tranciti, I built serverless microservices for a GPS geolocation platform, automated deployments, ran the 24/7 emergency operations line, and upgraded a route-monitoring mobile app to Android 12+. What sets me apart is sysadmin-level Linux: it is my daily driver, I am familiar with many distributions, and I have maintained a Gentoo package repository. I am looking to join a DevOps and Cloud team.

% =====================================================================
\section{Key Achievements}
% =====================================================================
\begin{itemize}
  \item Cut deployment time by 40\% through Jenkins automation and reduced production errors by 25\% with Docker and code quality tooling.
  \item Upgraded a real-time route-monitoring app to Android 12+ despite obsolete dependencies and sparse documentation.
  \item 24/7 emergency line (``red telephone''): resolved an early-morning outage of the real-time vehicle reporting service and trained the on-call rotation on the protocol.
  \item Built a WordPress site from scratch and had it live in under a week (Serakin), with 30\%+ higher user retention, and cut page load time by 30\% on the Kimún platform.
\end{itemize}

% =====================================================================
\section{Experience}
% =====================================================================

\cvjob{Junior Software Engineer}{Tranciti (Waypoint Telecomunicaciones)}{Oct 2021 -- Jan 2025}
\begin{itemize}
  % TODO: validate the 30% efficiency figure (the DOCX had "[validate]")
  \item Developed microservices and serverless functions with AWS Lambda (Java and Node.js) for a GPS geolocation platform, including an API to communicate with GPS devices; improved efficiency by 30\%.
  \item Implemented infrastructure as code with AWS CloudFormation and automated deployments with Jenkins, reducing deployment time by 40\%. Integrated Docker and code quality tools, which lowered production errors by 25\%.
  \item Upgraded the real-time route-monitoring mobile app to the latest Android SDK (Android 12+). Overcame obsolete dependencies, constant framework updates, and sparse documentation, and verified compatibility through testing on a company device.
  \item Managed the ``red telephone'' (24/7 emergency line). When the vehicle reporting service went down at 6 a.m., I reviewed the protocol and logs, restarted the service, confirmed the fix with the client, and proposed changes to prevent a recurrence. Trained the rotation on the protocol and on how to classify an emergency.
  \item Was among the team's strongest Linux users. Wrote scripts to run queries across multiple databases and optimized PostgreSQL and DynamoDB/Athena queries, with up to 20\% lower response times.
  \item Developed a frontend module in Angular for clients and supported mobile apps (Titanium, Ionic, Expo Go) for 5 clients.
\end{itemize}

\cvjob{Frontend Developer}{Innovación y Proyectos Kimun Ltda., Temuco}{Dec 2019}
\begin{itemize}
  \item Redesigned components of the ``Mundo de Vegetales'' site, such as the navigation bar, applying color psychology in CSS. Improved usability, increased retention by 20\%, and helped establish the site's design framework.
  \item Optimized site loading by removing unnecessary elements and refactoring code to reduce load cycles: 30\% shorter load time.
  \item Analyzed the codebase and delivered reports on its status and recommended procedures, discussed in daily stand-ups.
\end{itemize}

\cvjob{WordPress Developer}{Serakin, Viña del Mar}{Oct 2016}
\begin{itemize}
  \item Built from scratch, backend and frontend, the website for a physiotherapy center with no web presence: installed it on the hosting with WordPress and had it live in under a week.
  \item Designed the site around a defined pattern, creating widgets, installing plugins, and building a custom child theme: 30\% higher user retention.
  \item Trained the client's team to manage the site once it was finished.
\end{itemize}

\vspace{5pt}\textbf{Earlier Web Development Experience} \hfill 2011 -- 2015\par\vspace{1pt}
\begin{itemize}
  \item \textbf{Web Developer and IT Team Lead} | Del Puerto Comunicaciones | Jun 2011 -- Mar 2014: the IT area was created when I joined, and I ended up leading it; developed and maintained 5+ client websites (WordPress, HTML, CSS, jQuery).
  \item \textbf{Web Developer} | Oxalix | Mar 2013 -- Jan 2014: forms and modules with Python and Django (MySQL), using a more robust and scalable architecture.
  \item \textbf{Web Developer} | Pontificia Universidad Católica de Valparaíso | Sep 2014 -- Apr 2015: English-learning system (CONICYT project) with PHP, jQuery, and MySQL; improved user experience by 20\%.
  \item \textbf{Frontend Developer} | Everis | Oct -- Dec 2015: client dashboard for the Consorcio website with Bootstrap and jQuery.
\end{itemize}

% =====================================================================
\section{Education}
% =====================================================================
\textbf{B.Eng. in Computer Science (Ingeniería de Ejecución en Informática)} | Universidad Técnica Federico Santa María | 2012 -- 2015

% =====================================================================
\section{Key Skills}
% =====================================================================
\begin{itemize}
  \cvskill{Linux and Systems}{Sysadmin-level Linux (multiple distributions, package repository maintenance)\sep CLI tools\sep Git\sep GitHub\sep GitLab}
  \cvskill{Cloud and DevOps}{AWS (Lambda, EC2, CloudFormation, CloudWatch, CodeCommit, DynamoDB, Athena, Glue, Redshift, QuickSight)\sep Docker\sep Jenkins\sep CI/CD\sep SonarQube}
  \cvskill{Languages}{Java\sep JavaScript (Node.js)\sep Python\sep PHP\sep C/C++\sep SQL\sep TypeScript (intermediate)}
  \cvskill{Frameworks}{Angular\sep Vue.js\sep Quarkus\sep Django\sep Ionic\sep Expo Go\sep Bootstrap\sep Tailwind CSS}
  \cvskill{Data}{PostgreSQL\sep MySQL\sep ETL\sep web scraping (BeautifulSoup)}
  \cvskill{Data and AI}{Ollama\sep LLMs (Qwen, GPT-oss, Llama, Deepseek)\sep Tools and harnesses (Deepseek Harness, Opencode)}
  \cvskill{Methodologies}{CI/CD\sep automated QA\sep TDD\sep DDD\sep version control\sep AI-assisted development}
  \cvskill{Differentiator}{Sysadmin-level Linux + AWS serverless}
  \cvskill{Languages (spoken)}{Spanish (native)\sep English (advanced written, intermediate spoken)}
\end{itemize}

\end{document}
